% <Note title>: <one line on what the note establishes>.
% Build with assets/build_note.sh of the summary-notes skill.
\documentclass[11pt]{article}

\usepackage[margin=1in]{geometry}
\usepackage{amsmath}
\usepackage{amssymb}
\usepackage{graphicx}
\usepackage{booktabs}
\usepackage{xcolor}
\usepackage[font=small,labelfont=bf,labelsep=period]{caption}
\usepackage[hidelinks]{hyperref}

% Symbols, one macro per object.
\newcommand{\hMpc}{\,h\,{\rm Mpc}^{-1}}
\newcommand{\T}{^{\mathsf{T}}}

% =====================================================================
% Numbers
% =====================================================================
% Every number the text quotes is defined here, with the file (and key)
% it comes from. Part 1: results, refreshed when the analysis is rerun;
% set \numbersfinalfalse to print them in red while they are stale.
% Part 2: inputs, which a rerun does not change.
%
% Statements of the text that rest on Part 1 (re-check after a refresh):
%   - <statement> (section "<title>").
\newif\ifnumbersfinal
\numbersfinaltrue
\newcommand{\refresh}[1]{\ifnumbersfinal#1\else\textcolor{red}{#1}\fi}
% Part 1: results.
% <what> (<file>, <key>)
\newcommand{\ExampleResult}{\refresh{0.30}}
% Part 2: inputs.
\newcommand{\ExampleInput}{0.05}

\begin{document}

\begin{center}
  {\LARGE <Title>\par}
  \vspace{0.6em}
  {\large <Subtitle: the question in one line>\par}
  \vspace{0.6em}
  {\large <Month Year>\par}
  \ifnumbersfinal\else
    \vspace{0.6em}
    {\color{red}Numbers in red are being refreshed.\par}
  \fi
\end{center}
\vspace{1em}

\section{Summary}
\label{sec:summary}

\paragraph{The question.}
<One or two sentences.>

\paragraph{The answer.}
<One or two sentences with the headline numbers, through macros.>

\paragraph{Why.}
<The reasoning in one paragraph, pointing to sections by title.>

\paragraph{Limits.}
<What the result does not cover.>

\section{Setup and definitions}
\label{sec:setup}

\paragraph{Data.} <...>
\paragraph{Model and fit.} <...>
\paragraph{The statistic.} <Its equation, once.>

\section{<First claim, as a title>}
\label{sec:first}

<The claim as the first sentence. How to read the figure. What it means.>

\begin{figure}[tbp]
  \centering
  \includegraphics[width=3.5in]{plots/<name>/<name>.pdf}
  \caption{<What is drawn. The encodings. What is omitted.>}
  \label{fig:<name>}
\end{figure}

\appendix

\section{Method details}
\label{app:method}

\section{Numerical checks}
\label{app:checks}

\section{Where the code and results are}
\label{app:code}

\begin{center}
  \begin{tabular}{ll}
    \toprule
    result & script $\rightarrow$ output \\
    \midrule
    <...> & \texttt{<...>} \\
    \bottomrule
  \end{tabular}
\end{center}

\section{Open issues}
\label{app:open}

\begin{itemize}
  \item <One to three sentences each.>
\end{itemize}

\end{document}
